\documentclass[11pt,a4paper]{article}

% ====================================================================
% Informe ejecutivo cientifico - Capa 1 del TFM
% Banco de pruebas comparativo de detectores de regimen de mercado
% Compilar con: tectonic report/informe_capa1.tex  (auto-ejecuta bibtex)
% ====================================================================

\usepackage[T1]{fontenc}
\usepackage[utf8]{inputenc}
\usepackage[spanish,es-noshorthands]{babel}
\usepackage[a4paper,margin=2.4cm]{geometry}
\usepackage{booktabs}
\usepackage{array}
\usepackage{graphicx}
\usepackage{float}
\usepackage{xcolor}
\usepackage{caption}
\usepackage{amsmath,amssymb}
\usepackage{enumitem}
\usepackage{hyperref}

\graphicspath{{../results/}}
\captionsetup{font=small,labelfont=bf}
\hypersetup{colorlinks=true,linkcolor=black!70,citecolor=black!70,urlcolor=blue!50!black}

% Colores de la paleta usada en las figuras
\definecolor{ventlarga}{HTML}{2B5D8A}
\definecolor{ventcorta}{HTML}{C98A2B}
\definecolor{negcolor}{HTML}{7A7A7A}

\newcommand{\code}[1]{\texttt{\detokenize{#1}}}
\newcommand{\Dt}[1]{\textbf{#1}}
\setlength{\emergencystretch}{3em}  % evita overfull hbox por rutas \code no divisibles

\title{\textbf{Detecci\'on de reg\'imenes de mercado:\\ un marco de evaluaci\'on causal y comparativo}\\[2mm]
\large Primera capa del TFM --- banco de pruebas de 12 detectores}
\author{Oscar \\ \small M\'aster en Inteligencia Artificial aplicada a los Mercados Financieros (MIAX)}
\date{Junio de 2026}

\begin{document}
\maketitle

% ====================================================================
\begin{abstract}
\noindent Este informe sintetiza la primera capa de un TFM mayor sobre detecci\'on de
reg\'imenes de mercado. Su objeto \textbf{no} es proponer un detector concreto, sino
construir un \textbf{marco de evaluaci\'on causal, comparable y honesto} con el que juzgar
a muchos detectores bajo el mismo protocolo. Se implementan y eval\'uan \textbf{12
detectores} de \textbf{7 familias} (reglas, clustering, HMM, Markov-Switching,
GARCH/RS-GARCH, change-point y redes) con un protocolo \emph{walk-forward}
out-of-sample sin \emph{look-ahead}, sobre el S\&P~500 y un universo multiactivo
(1985--2026). El resultado central es que \textbf{no existe un detector dominante}: cuatro
familias se reparten seis ejes de calidad. La cobertura de crisis sist\'emicas grandes la
lideran el Markov-Switching de varianza (D5), el GARCH-t (D6) y una simple regla sobre el
VIX (D1); la especificidad, la persistencia, la anticipaci\'on y el coste los domina el
\emph{change-point} CUSUM (D7); y el ajuste estad\'istico (BIC) lo gana el HMM t-Student de
cuatro estados (D8). Adem\'as se documentan seis hallazgos metodol\'ogicos transversales
---el coraz\'on del trabajo--- y se recomienda, de forma \emph{consistente con} (no
``confirmando'') la propuesta del TFM mayor, un n\'ucleo HMM t-Student multi-estado
complementado con un change-point r\'apido. El informe es autocontenido y declara con
transparencia sus limitaciones: $n\approx4$ crisis sin contraste de significancia, BIC
in-sample, y lead/lag censurado por la ventana de b\'usqueda.
\end{abstract}

% ====================================================================
\section{Resumen ejecutivo}

\textbf{Qu\'e se hizo.} Se dise\~n\'o un marco \'unico de evaluaci\'on
(\code{src/evaluation.py}) y una interfaz com\'un de detector
(\code{src/detector_base.py}) con etiquetas can\'onicas ($0=$ calma $\dots$ $n{-}1=$
crisis). Cada uno de los 12 detectores se reentrena y predice en ventanas
\emph{walk-forward} expandidas, produciendo etiquetas \emph{out-of-sample} (OOS) sin ver el
futuro. Sobre esas etiquetas se calculan las mismas m\'etricas para todos: cobertura de
crisis, especificidad en ventanas-trampa, lead/lag, frecuencia de conmutaci\'on,
persistencia y ---donde el modelo lo permite--- bondad de ajuste (logL/AIC/BIC).

\textbf{Qu\'e se hall\'o.} (i)~Ninguna familia gana en todo: hay un \emph{mejor-para-qu\'e},
no un mejor absoluto. (ii)~Para cubrir crisis grandes y lentas basta con reaccionar a la
volatilidad con suficiente hist\'orico: el MS de varianza (D5, $0.98$), el GARCH-t (D6,
$0.97$) y la regla VIX (D1, $0.92$) empatan pr\'acticamente en cobertura de
GFC$+$COVID ---primer aviso de parsimonia. (iii)~El change-point CUSUM (D7) es el \'unico
con especificidad $1.00$ en las dos trampas, la mayor persistencia ($\approx436$ d\'ias) y
el menor coste. (iv)~El HMM t-Student de 4 estados (D8) ajusta mucho mejor que su gemelo
gaussiano ($\Delta\mathrm{BIC}\approx{+}10\,963$). (v)~Los dos modelos m\'as complejos
(MS-GARCH y autoencoder) \emph{no} superan a los baselines: la parsimonia queda validada.
(vi)~El \emph{taper tantrum} de 2013 es un punto ciego universal: ning\'un detector basado
en volatilidad lo ve, porque fue un shock de tipos sin estr\'es sist\'emico de renta
variable.

\textbf{Qu\'e se recomienda.} Llevar a la siguiente capa un \textbf{n\'ucleo HMM t-Student
multi-estado} (D8), \emph{consistente con} la propuesta original del TFM, emparejado con un
\textbf{change-point r\'apido} tipo D7 como sistema de alerta temprana / ``segunda
velocidad'', y conservar D1/D5/D6 como controles \emph{baseline}. Se descartan como
detectores operativos los exploratorios negativos D11 y D12.

% ====================================================================
\section{Introducci\'on y objetivo de la capa}

\subsection{De d\'onde viene este trabajo}
Esta capa nace del detector de reg\'imenes de una pr\'actica previa: un \textbf{HMM
gaussiano de 2 estados, ajustado in-sample y con \emph{look-ahead}}. Aquel modelo acertaba
las crisis sist\'emicas grandes (Lehman 2008: $98.6\%$ de d\'ias en crisis; COVID 2020:
$92.3\%$) pero se perd\'ia las correcciones r\'apidas (Taper Tantrum 2013: $10.9\%$;
sell-off del Q4 2018: $20.6\%$). De su an\'alisis cr\'itico salieron ocho limitaciones que
motivan este banco de pruebas: usar s\'olo 2 estados; ser ciego a crisis r\'apidas;
etiquetar por umbral \emph{post-hoc}; entrenar in-sample sin validaci\'on
\emph{walk-forward}; asumir emisiones gaussianas (que subestiman las colas); estandarizar
con estad\'isticos de toda la muestra (otra fuga de futuro sutil); una inicializaci\'on
degenerada; y no reportar incertidumbre (Viterbi ``duro'' en vez de probabilidades
suaves).

\subsection{Qu\'e propone el TFM mayor}
El TFM mayor define el sistema objetivo: un \textbf{HMM t-Student de 4 estados}, con
selecci\'on de $K$ por BIC y una l\'ogica de \textbf{``dos velocidades''} que distingue la
correcci\'on r\'apida de la crisis sist\'emica lenta. La pregunta que esta capa debe
responder con evidencia es: \emph{\textquestiondown la exploraci\'on de detectores respalda esa elecci\'on,
o sugiere ajustarla?}

\subsection{Qu\'e es ---y qu\'e no es--- esta capa}
Esta capa \textbf{no es un detector concreto}. Es un \textbf{marco de evaluaci\'on causal,
comparativo y honesto} que enfrenta 12 detectores de 7 familias bajo una sola interfaz y un
protocolo OOS com\'un. Su tesis central no es ``cu\'al es el mejor detector'' sino
\textbf{``cu\'al es el mejor para qu\'e''}: cada familia domina un eje distinto, y la
comparaci\'on s\'olo es leg\'itima si se hace \textbf{causal} y \textbf{con equidad de
ventana}.

% ====================================================================
\section{Datos y marco causal de evaluaci\'on}

\subsection{Datos}
Se descarga un universo multiactivo sin imputar (cada serie arranca en su fecha real):
S\&P~500, VIX, MOVE, treasuries (TLT, IEF), cr\'edito \emph{high yield} (HYG), oro (GLD),
d\'olar (DXY) y la pendiente de la curva ($\code{^TNX}-\code{^IRX}$). La ventana com\'un
sin huecos es \textbf{2007-04-11 $\to$ 2026-06-12} (gobernada por el inicio de HYG), de la
que se derivan \textbf{15 \emph{features} causales}. Varios detectores que s\'olo necesitan
S\&P~500 y/o VIX usan hist\'oricos m\'as largos (desde 1985/1990/1993), lo que les permite
evaluar la GFC de 2008 \emph{out-of-sample}.

Dos hechos estilizados del an\'alisis exploratorio gobiernan el dise\~no: las
\textbf{colas pesadas} (curtosis en exceso de $25.6$ en el S\&P~500 y $39.6$ en HYG;
Figura~\ref{fig:fattails}), que motivan emisiones t-Student y GARCH; y la \textbf{subida de
correlaciones en crisis} (Figura~\ref{fig:corr}), que motiva los detectores multivariantes
de turbulencia. La verdad-terreno de eventos (ventanas de crisis y trampas, y los suelos de
\emph{drawdown}) se fija sobre la serie real del S\&P~500 (Figura~\ref{fig:drawdown}).

\begin{figure}[H]
  \centering
  \includegraphics[width=0.92\textwidth]{eda_sp500_drawdown.png}
  \caption{S\&P~500 (escala log) y su \emph{drawdown} desde 1985, con las ventanas de
  \textbf{crisis} (rojo) y de \textbf{trampa / falso positivo} (naranja) que definen el
  protocolo de evaluaci\'on. Estas ventanas son la verdad-terreno frente a la que se mide
  cobertura, especificidad y lead/lag.}
  \label{fig:drawdown}
\end{figure}

\begin{figure}[H]
  \centering
  \begin{minipage}[t]{0.49\textwidth}
    \centering
    \includegraphics[width=\textwidth]{eda_fat_tails.png}
    \caption{Distribuci\'on de retornos diarios frente a la Normal (eje log). El exceso de
    curtosis ($25.6$ en el S\&P~500, $39.6$ en HYG) justifica colas pesadas y modelos
    t-Student/GARCH.}
    \label{fig:fattails}
  \end{minipage}\hfill
  \begin{minipage}[t]{0.49\textwidth}
    \centering
    \includegraphics[width=\textwidth]{eda_corr.png}
    \caption{Matriz de correlaci\'on de retornos diarios. La estructura
    renta-fija/renta-variable justifica el universo multiactivo de los detectores de
    turbulencia y clustering.}
    \label{fig:corr}
  \end{minipage}
\end{figure}

\subsection{El contrato de comparabilidad}
Tres principios, en la l\'inea metodol\'ogica del \emph{machine learning} financiero
causal~\cite{lopezdeprado2018}, hacen que la comparaci\'on sea leg\'itima:
\begin{description}[leftmargin=1.2em,style=nextline]
  \item[\emph{Features} causales.] Toda estandarizaci\'on es un \emph{z-score} expandido
  (en $t$ s\'olo se usan estad\'isticos de datos $\le t$). Nunca media/desviaci\'on de toda
  la muestra: \emph{eso} era el \emph{look-ahead} detectado en la tarea previa.
  \item[\emph{Walk-forward} OOS.] Nada se juzga in-sample. El detector se reentrena en
  ventanas expandidas y predice el bloque siguiente con par\'ametros congelados; cada fecha
  recibe la predicci\'on causal del fold que la cubri\'o por primera vez.
  \item[Misma interfaz, mismas m\'etricas.] Cada detector implementa \code{RegimeDetector}
  y se punt\'ua con \code{evaluation.py}. Las etiquetas se canonicalizan por criterio
  econ\'omico \textbf{vol-primario} (los estados se ordenan por su volatilidad; el retorno
  medio s\'olo desempata entre vol\'umenes pr\'oximos), y el orden se re-fija con los
  retornos del S\&P~500 \textbf{en cada fold}.
\end{description}

\noindent Este \'ultimo punto (la canonicalizaci\'on econ\'omica robusta por fold) es lo
que hace \emph{comparables} entre s\'i a familias que ni siquiera operan sobre retornos
(varianza, $\sigma$ GARCH, change-point, distancia de Mahalanobis). Sin \'el, la mitad de
los detectores habr\'ian sido incomparables o se habr\'ian invertido aleatoriamente.

\paragraph{M\'etricas.} (a)~\emph{Cobertura de crisis}: \% de d\'ias en ``crisis'' dentro
de 2008/2011/2020/2022 (alto = sensible). (b)~\emph{Especificidad}: no disparar de forma
sostenida en las trampas 2013/2018 (bajo = espec\'ifico). (c)~\emph{Lead/lag}: d\'ias entre
la se\~nal sostenida y el suelo del \emph{drawdown} (negativo = anticipa). (d)~\emph{Tasa de
falsas alarmas} global. (e)~\emph{Conmutaci\'on / persistencia} (penaliza el
\emph{flickering}). (f)~\emph{logL/AIC/BIC} donde el modelo lo permite.

% ====================================================================
\section{Los 12 detectores}

Se cubren las 7 familias, de \emph{baseline} a avanzado, m\'as dos exploratorios cuyo
resultado negativo es informativo (Tabla~\ref{tab:mapa}).

\begin{table}[H]
  \centering
  \small
  \caption{Mapa de los 12 detectores: familia, clase y ventana de evaluaci\'on OOS. La
  columna \emph{vio 2008 OOS} es decisiva: s\'olo los seis de hist\'orico largo evaluaron
  la GFC fuera de muestra.}
  \label{tab:mapa}
  \begin{tabular}{llll c}
    \toprule
    ID & Detector & Familia & Clase & Vio 2008 OOS \\
    \midrule
    D1 & \code{rule_vix_threshold} & F1 reglas & baseline & s\'i \\
    D2 & \code{rule_composite_riskoff} & F1 reglas & baseline & no \\
    D3 & \code{clustering_gmm_k3} & F2 clustering & baseline & no \\
    D4 & \code{hmm_gaussian_2s} & F3 HMM & baseline (puente) & no \\
    D5 & \code{markov_switching_var_2s} & F4 Markov-Switching & avanzado & s\'i \\
    D6 & \code{garch_t_vol} & F5 GARCH & avanzado & s\'i \\
    D7 & \code{changepoint_online} & F6 change-point & avanzado & s\'i \\
    D8 & \code{hmm_tstudent_4s} & F3 HMM (avanzado) & avanzado & no \\
    D9 & \code{jump_model} & F2$\leftrightarrow$F3 jump model & avanzado & no \\
    D10 & \code{turbulence_mahalanobis} & F1 multivariante & avanzado & s\'i \\
    D11 & \code{msgarch_regime} & F5 MS-GARCH & \textit{explor.-negativo} & s\'i \\
    D12 & \code{deep_ae_regime} & F7 redes & \textit{explor.-negativo} & no \\
    \bottomrule
  \end{tabular}
\end{table}

\begin{description}[leftmargin=0pt,labelsep=0.5em,itemsep=2pt]
  \item[D1 -- regla VIX.] \cite{reglas_bloom2009} Aut\'omata de 2 estados sobre el nivel del VIX con hist\'eresis y
  \emph{dwell}. Baseline de ``miedo'' limpio y persistente; entre los mejores en cobertura
  sist\'emica de ventana larga.
  \item[D2 -- regla compuesta.] Voto VIX + cr\'edito + curva + \emph{drawdown}. Mejora a D1
  en 2022 pero a costa de m\'as ruido; complementa, no sustituye.
  \item[D3 -- GMM clustering.] \cite{clust_pedregosa2011sklearn,clust_schwarz1978bic} Mezcla gaussiana $K{=}3$ sin din\'amica temporal. Buena
  cobertura pero \emph{flickering} severo: aisla cu\'anto aporta la din\'amica temporal.
  \item[D4 -- HMM gaussiano (puente).] \cite{hmm_rabiner1989,hmm_viterbi1967} R\'eplica causal del detector previo. Sirve para
  medir el efecto del \emph{look-ahead}; su supuesto gaussiano subestima las colas y motiva
  a D8.
  \item[D5 -- Markov-Switching de varianza.] \cite{hamilton1989,ms_kim1994} Probabilidades filtradas causales. Ganador de
  cobertura sist\'emica en ventana larga; control interpretable, pero caro.
  \item[D6 -- GARCH-t.] \cite{vol_bollerslev1986,vol_glosten1993} GJR-GARCH(1,1)-t con umbral causal sobre $\sigma$. Capta 2018 que
  D4 no; sensor de volatilidad fuerte y barato.
  \item[D7 -- change-point CUSUM.] \cite{cp_page1954,cp_truong2020} CUSUM de Page online robusto sobre la volatilidad.
  Domina cuatro ejes (especificidad, persistencia, lead/lag, coste); mejor candidato a
  ``segunda velocidad''.
  \item[D8 -- HMM t-Student 4 estados.] Colas pesadas con $\nu$ por estado. Gana el eje
  BIC; n\'ucleo recomendado, con matices (Secci\'on~\ref{sec:reco}).
  \item[D9 -- jump model.] \cite{hmm_nystrup2020} Modelo de saltos penalizados ($\lambda{=}50$). Anti-flickering
  rotundo, pero pierde cobertura de crisis lentas.
  \item[D10 -- turbulencia de Mahalanobis.] \cite{kritzman2010,kritzman2012} Distancia multivariante con covarianza
  expandida. Capta sist\'emicas pero no 2013 (no fue turbulencia conjunta).
  \item[D11 -- MS-GARCH \textit{(negativo)}.] \cite{vol_haasmittnikpaolella2004} MS-GARCH-t de Haas-Mittnik-Paolella. Degenera
  en \emph{walk-forward}: el fold de la GFC colapsa a un r\'egimen
  ($\mathrm{cov\_GFC}=0\%$).
  \item[D12 -- autoencoder \textit{(negativo)}.] \cite{nn_hinton2006} AE $\to$ GMM frente a baseline PCA $\to$
  GMM. El autoencoder empeora al PCA sin ganar cobertura: la no linealidad no aporta con
  pocos datos.
\end{description}

% ====================================================================
\section{Resultados comparativos}
\label{sec:resultados}

\subsection{Tabla maestra}
La Tabla~\ref{tab:maestra} resume el comportamiento global. La separaci\'on por ventana es
obligatoria: las dos mitades no son directamente comparables en cobertura.

\begin{table}[H]
  \centering
  \small
  \caption{Tabla maestra resumida (valores \emph{verbatim} de
  \code{results/metrics_master.csv}). \emph{switching}: conmutaciones por d\'ia
  (bajo = persistente); \emph{dur.}: duraci\'on media de r\'egimen en d\'ias; \emph{FA
  global}: tasa global de falsa alarma; BIC s\'olo en modelos generativos.}
  \label{tab:maestra}
  \begin{tabular}{ll c r r r r}
    \toprule
    ID & Detector & Vent. & switching & dur.\,(d) & FA global & BIC \\
    \midrule
    \multicolumn{7}{l}{\textit{\color{ventlarga}Ventana larga (vio 2008 OOS)}}\\
    D6 & \code{garch_t_vol} & 1993--2026 & 0.0141 & 70.15 & 0.845 & 26\,627 \\
    D5 & \code{markov_switching_var_2s} & 1993--2026 & 0.0557 & 17.92 & 0.774 & 28\,024 \\
    D7 & \code{changepoint_online} & 1993--2026 & 0.0022 & 435.68 & 0.867 & --- \\
    D1 & \code{rule_vix_threshold} & 1998--2026 & 0.0132 & 75.20 & 0.697 & --- \\
    D10 & \code{turbulence_mahalanobis} & 1998--2026 & 0.0873 & 11.44 & 0.815 & --- \\
    D11 & \code{msgarch_regime} \textit{(neg.)} & 1991--2026 & 0.0312 & 31.93 & 0.949 & 26\,823 \\
    \midrule
    \multicolumn{7}{l}{\textit{\color{ventcorta}Ventana corta (no vio 2008 OOS)}}\\
    D3 & \code{clustering_gmm_k3} & 2015--2026 & 0.1261 & 7.91 & 0.494 & 63\,016 \\
    D4 & \code{hmm_gaussian_2s} & 2012--2026 & 0.1004 & 9.93 & 0.731 & 35\,379 \\
    D2 & \code{rule_composite_riskoff} & 2015--2026 & 0.0385 & 25.72 & 0.725 & --- \\
    D8 & \code{hmm_tstudent_4s} & 2012--2026 & 0.0520 & 19.13 & 0.519 & 24\,416 \\
    D9 & \code{jump_model} & 2015--2026 & 0.0053 & 176.60 & 0.624 & --- \\
    D12 & \code{deep_ae_regime} \textit{(neg.)} & 2015--2026 & 0.2873 & 3.48 & 0.603 & --- \\
    \bottomrule
  \end{tabular}
\end{table}

\subsection{Cobertura de crisis: estricta y estr\'es agregado}
La cobertura ``estricta'' cuenta los d\'ias en el estado de crisis ($\code{state}=n{-}1$).
Para comparar con justicia los detectores multi-estado con los binarios se define adem\'as
el \textbf{estr\'es agregado}: en binarios coincide con la crisis; en multi-estado
($n\ge3$) es la uni\'on de los dos estados m\'as severos (correcci\'on $+$ crisis). Las
Tablas~\ref{tab:cobestricta} y~\ref{tab:estres} reportan ambas.

\begin{table}[H]
  \centering
  \small
  \caption{Cobertura de crisis con definici\'on \textbf{estricta} (proporci\'on de d\'ias
  en crisis por ventana). \emph{NaN} = ventana fuera del rango OOS del detector (no se
  penaliza lo que no se pudo ver).}
  \label{tab:cobestricta}
  \begin{tabular}{ll cccc}
    \toprule
    ID & Detector & GFC 2008 & EuroDeuda 2011 & COVID 2020 & Inflaci\'on 2022 \\
    \midrule
    \multicolumn{6}{l}{\textit{\color{ventlarga}Ventana larga}}\\
    D6 & \code{garch_t_vol} & 1.000 & 0.741 & 0.940 & 0.804 \\
    D5 & \code{markov_switching_var_2s} & 0.993 & 0.741 & 0.960 & 0.737 \\
    D7 & \code{changepoint_online} & 1.000 & 0.671 & 0.840 & 0.766 \\
    D1 & \code{rule_vix_threshold} & 0.938 & 0.635 & 0.900 & 0.349 \\
    D10 & \code{turbulence_mahalanobis} & 0.822 & 0.482 & 0.760 & 0.431 \\
    D11 & \code{msgarch_regime} \textit{(neg.)} & 0.000 & 0.118 & 0.200 & 0.359 \\
    \midrule
    \multicolumn{6}{l}{\textit{\color{ventcorta}Ventana corta} (GFC y EuroDeuda fuera de OOS)}\\
    D3 & \code{clustering_gmm_k3} & NaN & NaN & 0.960 & 0.865 \\
    D4 & \code{hmm_gaussian_2s} & NaN & NaN & 0.960 & 0.861 \\
    D2 & \code{rule_composite_riskoff} & NaN & NaN & 0.840 & 0.538 \\
    D8 & \code{hmm_tstudent_4s} & NaN & NaN & 0.660 & 0.332 \\
    D9 & \code{jump_model} & NaN & NaN & 0.720 & 0.168 \\
    D12 & \code{deep_ae_regime} \textit{(neg.)} & NaN & NaN & 0.540 & 0.101 \\
    \bottomrule
  \end{tabular}
\end{table}

\begin{table}[H]
  \centering
  \small
  \caption{Cobertura con \textbf{estr\'es agregado} (correcci\'on $+$ crisis) para los tres
  multi-estado; en los binarios coincide con la estricta. El criterio relajado eleva mucho
  la cobertura de crisis lentas de D8 y D12.}
  \label{tab:estres}
  \begin{tabular}{ll cc cc}
    \toprule
    & & \multicolumn{2}{c}{COVID 2020} & \multicolumn{2}{c}{Inflaci\'on 2022} \\
    \cmidrule(lr){3-4}\cmidrule(lr){5-6}
    ID & Detector ($K$) & estricta & estr\'es & estricta & estr\'es \\
    \midrule
    D3 & \code{clustering_gmm_k3} ($K{=}3$) & 0.960 & 0.960 & 0.865 & 0.865 \\
    D8 & \code{hmm_tstudent_4s} ($K{=}4$) & 0.660 & \textbf{0.960} & 0.332 & \textbf{0.904} \\
    D12 & \code{deep_ae_regime} ($K{=}3$) & 0.540 & \textbf{0.960} & 0.101 & \textbf{0.788} \\
    \bottomrule
  \end{tabular}
\end{table}

\noindent El estr\'es agregado tiene \textbf{dos caras}, y mostrar ambas es la lectura
honesta (Figura~\ref{fig:estres}): elevar a ``correcci\'on $+$ crisis'' sube la cobertura
de crisis lentas, \emph{pero tambi\'en} la activaci\'on en las trampas. En D8 el sell-off
de 2018 pasa de $0.034$ a $0.814$ y la tasa global de falsa alarma de $0.52$ a $0.79$; en
D3, la trampa de 2018 salta de $0.000$ a $0.729$ (Tabla~\ref{tab:especificidad}). No es una
mejora gratuita, es un \emph{trade-off}.

\begin{table}[H]
  \centering
  \small
  \caption{Especificidad en las ventanas-trampa (proporci\'on de d\'ias ``encendido'';
  bajo = mejor). Se muestran la versi\'on estricta y la de estr\'es agregado. \emph{NaN} =
  trampa fuera del OOS del detector.}
  \label{tab:especificidad}
  \begin{tabular}{ll cc cc}
    \toprule
    & & \multicolumn{2}{c}{Taper 2013} & \multicolumn{2}{c}{Sell-off 2018} \\
    \cmidrule(lr){3-4}\cmidrule(lr){5-6}
    ID & Detector & estricta & estr\'es & estricta & estr\'es \\
    \midrule
    D7 & \code{changepoint_online} & 0.000 & 0.000 & 0.000 & 0.000 \\
    D1 & \code{rule_vix_threshold} & 0.000 & 0.000 & 0.063 & 0.063 \\
    D5 & \code{markov_switching_var_2s} & 0.038 & 0.038 & 0.810 & 0.810 \\
    D6 & \code{garch_t_vol} & 0.113 & 0.113 & 0.873 & 0.873 \\
    D10 & \code{turbulence_mahalanobis} & 0.123 & 0.123 & 0.302 & 0.302 \\
    D9 & \code{jump_model} & NaN & NaN & 0.000 & 0.000 \\
    D3 & \code{clustering_gmm_k3} & NaN & NaN & 0.000 & \textbf{0.729} \\
    D8 & \code{hmm_tstudent_4s} & 0.000 & 0.000 & 0.034 & \textbf{0.814} \\
    D12 & \code{deep_ae_regime} \textit{(neg.)} & NaN & NaN & 0.153 & 0.186 \\
    D11 & \code{msgarch_regime} \textit{(neg.)} & 0.000 & 0.000 & 0.937 & 0.937 \\
    \bottomrule
  \end{tabular}
\end{table}

\subsection{Los planos de \emph{trade-off}}
Las Figuras~\ref{fig:sensesp} y~\ref{fig:persens} cruzan los ejes en lugar de apilar tablas.
Para que la comparaci\'on sea justa entre ventanas largas y cortas, la \textbf{sensibilidad
se mide en la ventana OOS com\'un a los doce (COVID 2020)}; la cobertura de la GFC se
ranquea s\'olo dentro del bloque de ventana larga.

\begin{figure}[H]
  \centering
  \begin{minipage}[t]{0.49\textwidth}
    \centering
    \includegraphics[width=\textwidth]{pdf_sensibilidad_especificidad.png}
    \caption{Plano sensibilidad $\leftrightarrow$ especificidad. Eje X: especificidad
    $=1-$ media de falsa alarma en 2013/2018. El color/forma codifica el grupo de ventana;
    D11/D12 son exploratorio-negativos.}
    \label{fig:sensesp}
  \end{minipage}\hfill
  \begin{minipage}[t]{0.49\textwidth}
    \centering
    \includegraphics[width=\textwidth]{pdf_persistencia_sensibilidad.png}
    \caption{Plano persistencia (duraci\'on media, eje log) $\leftrightarrow$ sensibilidad.
    D3/D12 \emph{flickean}; D7 y las reglas son muy persistentes.}
    \label{fig:persens}
  \end{minipage}
\end{figure}

\begin{figure}[H]
  \centering
  \begin{minipage}[t]{0.49\textwidth}
    \centering
    \includegraphics[width=\textwidth]{pdf_bic.png}
    \caption{BIC de los modelos generativos (menor = mejor). \textbf{Aviso}: s\'olo es
    estrictamente comparable sobre las mismas \emph{features} (D4 frente a D8). Gris =
    negativo.}
    \label{fig:bic}
  \end{minipage}\hfill
  \begin{minipage}[t]{0.49\textwidth}
    \centering
    \includegraphics[width=\textwidth]{pdf_estres_vs_estricta.png}
    \caption{Multi-estado: cobertura estricta frente a estr\'es agregado en COVID 2020 e
    Inflaci\'on 2022. Lectura honesta de los dos lados del \emph{trade-off}.}
    \label{fig:estres}
  \end{minipage}
\end{figure}

\begin{figure}[H]
  \centering
  \includegraphics[width=0.86\textwidth]{pdf_rank_heatmap.png}
  \caption{Ranking por eje (1 = mejor; m\'as oscuro = mejor rango). \textbf{Corrige el
  defecto} de la versi\'on previa: la cobertura se separa por ventana ---la GFC 2008
  s\'olo ranquea entre los seis detectores de ventana larga, COVID 2020 entre los doce--- y
  las filas se agrupan por grupo de ventana. El lead/lag medio est\'a censurado (ver
  Secci\'on~\ref{sec:limitaciones}).}
  \label{fig:rank}
\end{figure}

\begin{figure}[H]
  \centering
  \includegraphics[width=0.92\textwidth]{pdf_leadlag.png}
  \caption{Lead/lag por evento (d\'ias de \emph{trading}; negativo = la se\~nal sostenida
  anticipa el suelo del \emph{drawdown}). El asterisco y el borde discontinuo marcan los
  valores \textbf{censurados} en $\pm252$ d\'ias: la se\~nal ya estaba activa al inicio de
  la ventana de b\'usqueda, de modo que no es anticipaci\'on genuina sino sesgo
  \emph{always-on}. Filas separadas por grupo de ventana.}
  \label{fig:leadlag}
\end{figure}

\subsection{Reparto del podio}
La Tabla~\ref{tab:podio} resume qu\'e familia gana en qu\'e eje. \textbf{Cuatro familias se
reparten seis ejes; no existe detector dominante.}

\begin{table}[H]
  \centering
  \small
  \caption{Reparto del podio por eje. La cobertura sist\'emica se da separada por grupo de
  ventana (nunca mezclada).}
  \label{tab:podio}
  \begin{tabular}{lll}
    \toprule
    Eje & Familia ganadora & Detalle \\
    \midrule
    Cobertura sist\'emica (ventana larga) & F4 Markov-Switching & D5 0.98 $>$ D6 0.97 $\approx$ D1/D7 0.92 \\
    Cobertura sist\'emica (ventana corta) & F2/F3 baseline & D3 0.96 $=$ D4 0.96 $>$ D2 0.84 \\
    Especificidad (trampas) & F6 change-point & D7 $1.00$ en ambas trampas \\
    Persistencia / anti-flicker & F6 change-point & D7 436 d; switching 0.002 \\
    Lead/lag (sostenido) & F6 change-point & D7 anticipa temprano en todos los eventos \\
    Ajuste BIC (in-sample) & F3 HMM t-Student & D8 24\,416; $\Delta$BIC $+10\,963$ vs D4 \\
    Coste computacional & F6 / F1 & D7 / D1 (bajo) \\
    \bottomrule
  \end{tabular}
\end{table}

\subsection{Scorecard e incertidumbre de cobertura}
La Figura~\ref{fig:scorecard} condensa el reparto del podio en un \emph{scorecard}: seis
ejes normalizados a $[0,1]$ por columna para los doce detectores. La lectura es inmediata
---ning\'un detector es verde en todas las columnas: ``mejor-para-qu\'e'' hecho imagen.

\begin{figure}[H]
  \centering
  \includegraphics[width=0.84\textwidth]{synth_scorecard.png}
  \caption{Scorecard ``mejor-para-qu\'e'': seis ejes normalizados por columna (verde = mejor
  en ese eje). Ning\'un detector domina todas las columnas; cada familia brilla en su eje.}
  \label{fig:scorecard}
\end{figure}

Como respuesta directa al trabajo futuro B1 (Secci\'on~\ref{sec:limitaciones}), se a\~nade una
cuantificaci\'on de incertidumbre \emph{intra-evento}: para cada ventana de crisis se calcula
un intervalo de confianza al 95\,\% de la cobertura mediante \emph{bootstrap por bloques}
(bloques de 5 d\'ias, que preservan la autocorrelaci\'on de los episodios). La
Figura~\ref{fig:covci} lo muestra para COVID-2020, la \'unica ventana OOS com\'un a los doce
detectores. Las bandas son anchas ---coherente con $n\approx4$ crisis--- y dejan ver que
diferencias de pocos puntos de cobertura \emph{no} son distinguibles, en l\'inea con la
cautela declarada en las limitaciones.

\begin{figure}[H]
  \centering
  \includegraphics[width=0.80\textwidth]{synth_coverage_ci.png}
  \caption{Cobertura de COVID-2020 con intervalo de confianza al 95\,\% por \emph{bootstrap}
  de bloques. El intervalo es \emph{intra-evento} (muestreo de d\'ias autocorrelados), no
  entre eventos: no sustituye a un contraste de significancia con $n\approx4$ crisis.}
  \label{fig:covci}
\end{figure}

% ====================================================================
\section{Hallazgos metodol\'ogicos transversales}

Son el verdadero producto de esta capa: lecciones sobre \emph{c\'omo} evaluar reg\'imenes
de forma honesta, transferibles m\'as all\'a de cualquier detector concreto.

\paragraph{(a) El \emph{look-ahead} compraba suavidad, no acierto.} Al reimplementar el HMM
gaussiano de la tarea previa de forma causal (D4), la conmutaci\'on sube de $0.047$
(in-sample) a $0.100$ (causal) y el modelo sigue fallando 2013 ($25\%$) y 2018 ($46\%$). El
acierto en crisis grandes nunca depende del \emph{look-ahead}; lo que \'este aportaba era
\textbf{cosm\'etica temporal} (persistencia falsa ``prestada del futuro''), no capacidad de
clasificaci\'on.

\paragraph{(b) El Viterbi por bloque era menos causal y menos estable que el filtrado
forward.} Al sustituir el Viterbi intra-bloque por \textbf{filtrado forward} en D4/D8, la
conmutaci\'on baj\'o de $0.124$ a $0.100$ y la duraci\'on media subi\'o de $8.1$ a $9.9$
d\'ias. Contra la intuici\'on, el cambio mejor\'o la estabilidad: el Viterbi reiniciaba la
decodificaci\'on en cada frontera de bloque (21 d\'ias), inyectando conmutaci\'on
artificial; el filtro forward con \emph{burn-in} es continuo ---m\'as causal \emph{y} m\'as
persistente. El ``\'optimo'' offline introduce un artefacto de frontera que el filtrado
causal evita.

\paragraph{(c) La t-Student est\'a bien especificada; el gaussiano no.} Con las mismas 7
\emph{features}, el HMM t-Student de 4 estados (D8) bate al gaussiano de 2 estados (D4) por
$\Delta\mathrm{BIC}\approx{+}10\,963$ (Tabla~\ref{tab:bic}). M\'as a\'un, D8 estima los
grados de libertad $\nu$ \textbf{decrecientes por estado} $[10.2,\,7.6,\,4.2,\,2.4]$: la
calma es casi gaussiana y la crisis tiene las colas m\'as pesadas, exactamente la
estructura que predice la teor\'ia de los hechos estilizados~\cite{hmm_bulla2011,
hmm_nystrup2015}. La t-Student no es un adorno: es la correcci\'on distribucional que la
curtosis 25--40 de los datos exige.

\begin{table}[H]
  \centering
  \small
  \caption{Bondad de ajuste de los modelos generativos. El BIC s\'olo es comparable sobre
  las mismas \emph{features}: el contraste limpio es D4 frente a D8 (ambos HMM, mismo
  conjunto de 7 \emph{features} puente).}
  \label{tab:bic}
  \begin{tabular}{ll r r r}
    \toprule
    ID & Detector & logL & AIC & BIC \\
    \midrule
    D8 & \code{hmm_tstudent_4s} & $-11\,536.3$ & 23\,390.7 & \textbf{24\,415.9} \\
    D6 & \code{garch_t_vol} & $-13\,285.6$ & 26\,583.1 & 26\,626.6 \\
    D11 & \code{msgarch_regime} \textit{(neg.)} & $-13\,365.5$ & 26\,750.9 & 26\,823.3 \\
    D5 & \code{markov_switching_var_2s} & $-13\,984.2$ & 27\,980.3 & 28\,023.8 \\
    D4 & \code{hmm_gaussian_2s} & $-17\,381.4$ & 34\,908.7 & \textbf{35\,379.4} \\
    D3 & \code{clustering_gmm_k3} & $-29\,789.0$ & 60\,392.0 & 63\,016.2 \\
    \midrule
    \multicolumn{5}{l}{$\Delta\mathrm{BIC}$ (D4 $-$ D8) $= 35\,379.4 - 24\,415.9 = \mathbf{10\,963.5}$ (evidencia ``muy fuerte'')}\\
    \bottomrule
  \end{tabular}
\end{table}

\paragraph{(d) El \emph{taper tantrum} de 2013 es un punto ciego universal.} Seis
detectores independientes de familias distintas que s\'i tuvieron 2013 en su OOS coinciden
en \emph{no} verlo (activaci\'on entre $0$ y $\sim12\%$): D1 reglas ($0.00$), D5 MS-VAR
($0.038$), D6 GARCH ($0.113$), D7 change-point ($0.00$), D8 HMM-t ($0.00$) y D10 turbulencia
($0.123$). El hallazgo \textbf{no} es ``2013 es ruido reclasificable''; es que \textbf{ninguna
definici\'on de crisis basada en volatilidad o correlaci\'on de renta variable captura
2013}, porque el \emph{taper} fue un shock de \emph{tipos} sin estr\'es sist\'emico de
renta variable. Que seis detectores converjan es \emph{evidencia robusta}, no un fallo
com\'un: \textbf{la taxonom\'ia de r\'egimen importa}. Capturar 2013 exigir\'ia ampliar la
taxonom\'ia de \emph{features} (curva de tipos, MOVE) o de reg\'imenes, no afinar un umbral;
por eso 2013 se mantiene como ventana-trampa para medir especificidad.

\paragraph{(e) La complejidad no se paga: parsimonia validada.} Los dos modelos m\'as
complejos son resultados negativos expl\'icitos. El MS-GARCH-t (D11) \emph{degenera} en
\emph{walk-forward}: el fold de la GFC colapsa a un \'unico r\'egimen, con cobertura de
2008 del $0\%$ y tasa de falsa alarma del $0.95$. El autoencoder (D12) \emph{empeora} a su
baseline lineal PCA $\to$ GMM (conmutaci\'on $0.287$ frente a $0.091$, falsa alarma $0.60$
frente a $0.14$) sin ganar cobertura. Con $\sim4$ crisis en la muestra, la riqueza
param\'etrica no se identifica: a veces es neutra (D12), a veces destructiva (D11). Los
baselines cubren los huecos que los complejos dejan.

\paragraph{(f) Dos arreglos imprescindibles del marco de evaluaci\'on.} (i)~\textbf{Severidad
vol-primaria}: con $K{=}2$, ordenar los estados por $z(\text{std})-z(\text{media})$ dejaba
que el signo ruidoso de una diferencia de medias casi nula invirtiera crisis$\leftrightarrow$calma
en detectores que separan s\'olo en varianza ($\sigma$ GARCH, turbulencia). La soluci\'on
fue ordenar por banda de volatilidad (ancho $15\%$ de la vol media) y desempatar por retorno
s\'olo entre vol\'umenes pr\'oximos. (ii)~\textbf{Etiquetado econ\'omico por fold}: re-fijar
el orden $0=$ calma $\dots$ $n{-}1=$ crisis con los retornos de \emph{cada} fold. Sin estos
dos arreglos, la mitad de los detectores habr\'ian sido incomparables.

% ====================================================================
\section{Recomendaci\'on para la siguiente capa}
\label{sec:reco}

La propuesta del TFM mayor es un HMM t-Student de 4 estados con l\'ogica de ``dos
velocidades''. \textbf{La evidencia de esta capa es \emph{consistente con} ese n\'ucleo y no
lo contradice; sugiere complementarlo.} Se evita deliberadamente el verbo ``confirma'': el
respaldo a D8 descansa en un eje \textbf{in-sample} (BIC) y en una m\'etrica
\textbf{favorable por construcci\'on} a $K\ge3$ (el estr\'es agregado), \emph{no} en
cobertura OOS estricta, donde D8 es modesto.

\begin{enumerate}[leftmargin=1.4em]
  \item \textbf{N\'ucleo HMM t-Student multi-estado (D8), reforzado por dos v\'ias ---cada
  una con su asterisco.}
    \begin{itemize}
      \item \emph{Por BIC (in-sample)}: gana el eje de ajuste con holgura
      ($\Delta\mathrm{BIC}\approx{+}10\,963$ sobre D4 con las mismas \emph{features}).
      \emph{Pero} el BIC es bondad de ajuste in-sample, no capacidad predictiva OOS.
      \item \emph{Por la l\'ogica correcci\'on$\leftrightarrow$crisis}: con estr\'es
      agregado, D8 iguala a los binarios (COVID $0.66\to0.96$; Inflaci\'on $0.33\to0.90$),
      lo que da sentido al cuarto estado y a la ``segunda velocidad''. \emph{Pero} el
      estr\'es agregado favorece por construcci\'on a los multi-estado, la cobertura OOS
      estricta es modesta (COVID $0.66$) y D8 nunca evalu\'o 2008 OOS. Es un argumento de
      \emph{plausibilidad estructural}, no de superioridad OOS demostrada.
    \end{itemize}
  \item \textbf{Complementar con un change-point r\'apido (D7)} como segunda velocidad /
  alerta temprana: domina especificidad, persistencia, lead/lag y coste. Donde el HMM-t da
  la \emph{taxonom\'ia} (qu\'e tipo de r\'egimen), el change-point da la \emph{reactividad}
  (cu\'ando cambia el nivel).
  \item \textbf{Conservar baselines de control}: D1 (regla VIX) y D5/D6 (MS-VAR / GARCH-t),
  con la mejor cobertura sist\'emica en ventana larga. Si un modelo complejo no bate a D1 en
  cobertura de GFC$+$COVID, no se justifica su coste.
  \item \textbf{Descartar D11 y D12} como detectores operativos: validan la parsimonia, no
  la aportan.
\end{enumerate}

\noindent En s\'intesis: la exploraci\'on \emph{no obliga} a cambiar la propuesta, pero la
matiza. El cuarto estado tiene sentido (es la ``correcci\'on'' que separa de la ``crisis'')
y la t-Student es necesaria, pero la validaci\'on definitiva del n\'ucleo deber\'a hacerse
con cobertura OOS estricta sobre m\'as crisis y, a ser posible, con la GFC dentro de la
muestra de prueba.

% ====================================================================
\section{Limitaciones y trabajo futuro}
\label{sec:limitaciones}

La honestidad de este marco exige declarar sus l\'imites:

\begin{description}[leftmargin=1.2em,style=nextline]
  \item[Significancia con $n\approx4$ crisis.] Toda la comparaci\'on descansa en cuatro
  crisis sist\'emicas (2008, 2011, 2020, 2022) y dos ventanas-trampa (2013, 2018). Se a\~naden
  \textbf{intervalos de confianza intra-evento} de la cobertura por \emph{bootstrap} de bloques
  (Figura~\ref{fig:covci}), pero \textbf{sigue sin haber contraste de significancia entre
  eventos}; diferencias de pocos puntos entre detectores \emph{no} son estad\'isticamente
  distinguibles. El podio es evidencia
  cualitativa direccional, no una ordenaci\'on con respaldo inferencial. Es una limitaci\'on
  asumida del dise\~no (las crisis son raras por definici\'on), no un descuido.
  \item[BIC in-sample.] Las m\'etricas de comportamiento (cobertura, falsa alarma, lead/lag,
  persistencia) son OOS sin \emph{look-ahead}; pero logL/AIC/BIC son \textbf{in-sample por
  definici\'on}. El eje de BIC mide especificaci\'on, no capacidad predictiva fuera de
  muestra, y as\'i se interpreta.
  \item[Lead/lag censurado.] El lead/lag se busca en los 252 d\'ias previos al suelo; un
  valor de $-252$ \textbf{no} significa ``252 d\'ias exactos de anticipaci\'on'', sino que
  la se\~nal ya estaba activa al inicio de la ventana ---el lead real es $\ge252$ y queda
  sin medir (censura por la derecha). D7 toca ese tope en 3 de sus 4 eventos, as\'i que su
  anticipaci\'on es un l\'imite inferior (Figura~\ref{fig:leadlag}); su virtud es \emph{cualitativa} (cruza pronto y se
  mantiene), no una cifra precisa de d\'ias.
  \item[Equidad de ventana.] La GFC 2008 cae cerca del inicio de datos para las ventanas
  cortas: D8 nunca la evalu\'o OOS y su validaci\'on sist\'emica se apoya en COVID 2020. Por
  el mismo motivo, 2013 es N/A (fuera de OOS) para los detectores de ventana corta y no
  cuenta ni a favor ni en contra.
\end{description}

\paragraph{Trabajo futuro declarado.} (B1)~\textbf{Bootstrap por bloques}
(\emph{stationary/circular block bootstrap}) sobre las series OOS ---\emph{parcialmente
abordado}: ya se proveen bandas de confianza de la cobertura por bloques
(Figura~\ref{fig:covci}); queda extenderlo a falsa alarma y conmutaci\'on, y a contrastes
pareados entre detectores respetando la autocorrelaci\'on. (B2)~\textbf{An\'alisis de sensibilidad de
hiperpar\'ametros} del protocolo \emph{walk-forward} (\code{train_size}, \code{step}) y del
etiquetado econ\'omico (\code{VOL_CLOSE_FRAC}), para comprobar que el reparto del podio es
robusto y no un artefacto de la configuraci\'on elegida.

% ====================================================================
\section*{Nota de reproducibilidad}
\addcontentsline{toc}{section}{Nota de reproducibilidad}
Todos los n\'umeros de este informe provienen \emph{verbatim} de un \'unico master
can\'onico, \code{results/metrics_master.csv} (43 columnas: comportamiento, ajuste, silueta
e intervalos de cobertura), reconstruible con \code{_rebuild_master.py}. Las figuras
\code{pdf_*.png} y \code{synth_*.png} (scorecard e IC de cobertura) se regeneran de forma
determinista desde ese mismo CSV con los scripts de \code{scripts/figs_pdf/}; las
\code{eda_*.png} provienen del an\'alisis exploratorio. Las m\'etricas de comportamiento y de
ajuste son id\'enticas a la versi\'on previa: esta revisi\'on solo \textbf{a\~nade} dos
m\'etricas ---silueta de reg\'imenes e intervalos de cobertura por \emph{bootstrap}--- sin
re-tunear ning\'un detector ni alterar ninguna cifra anterior.

% ====================================================================
\bibliographystyle{plain}
\bibliography{references}

\end{document}
